\documentclass[10pt,a4paper]{amsart}
\usepackage{amsmath,amssymb,amsthm,graphicx}
% Retained minimal preamble for the Stacks Project chapter more-algebra.tex (More on Algebra).
% The chapter file pinned-original-source.tex carries only \input{preamble} on its first line, so no
% preamble text lives inside the chapter itself.  The Stacks Project uses one single shared preamble for
% every chapter; the snapshot stages it identically as sources/stacks-chow/preamble.tex and
% sources/stacks-etale-cohomology/preamble.tex.  The block below is copied mechanically, line for line,
% from sources/stacks-chow/preamble.tex lines 14-18 (the Xy-pic setup) and 154-217 (the native theorem
% declarations together with the mathematical macros).  Nothing was added, removed or redefined; the
% wrapper class and the amsmath/amssymb/amsthm/graphicx base packages are supplied by the assembler.
% Chapter: more-algebra.tex -- More on Algebra.
% For commutative diagrams we use Xy-pic
\usepackage[all]{xy}

% We use 2cell for 2-commutative diagrams.
\xyoption{2cell}

% Theorem environments.
%
\theoremstyle{plain}
\newtheorem{theorem}[subsection]{Theorem}
\newtheorem{proposition}[subsection]{Proposition}
\newtheorem{lemma}[subsection]{Lemma}

\theoremstyle{definition}
\newtheorem{definition}[subsection]{Definition}
\newtheorem{example}[subsection]{Example}
\newtheorem{exercise}[subsection]{Exercise}
\newtheorem{situation}[subsection]{Situation}

\theoremstyle{remark}
\newtheorem{remark}[subsection]{Remark}
\newtheorem{remarks}[subsection]{Remarks}

\numberwithin{equation}{subsection}

% Macros
%
\def\lim{\mathop{\mathrm{lim}}\nolimits}
\def\colim{\mathop{\mathrm{colim}}\nolimits}
\def\Spec{\mathop{\mathrm{Spec}}}
\def\Hom{\mathop{\mathrm{Hom}}\nolimits}
\def\Ext{\mathop{\mathrm{Ext}}\nolimits}
\def\SheafHom{\mathop{\mathcal{H}\!\mathit{om}}\nolimits}
\def\SheafExt{\mathop{\mathcal{E}\!\mathit{xt}}\nolimits}
\def\Sch{\mathit{Sch}}
\def\Mor{\mathop{\mathrm{Mor}}\nolimits}
\def\Ob{\mathop{\mathrm{Ob}}\nolimits}
\def\Sh{\mathop{\mathit{Sh}}\nolimits}
\def\NL{\mathop{N\!L}\nolimits}
\def\CH{\mathop{\mathrm{CH}}\nolimits}
\def\proetale{{pro\text{-}\acute{e}tale}}
\def\etale{{\acute{e}tale}}
\def\QCoh{\mathit{QCoh}}
\def\Ker{\mathop{\mathrm{Ker}}}
\def\Im{\mathop{\mathrm{Im}}}
\def\Coker{\mathop{\mathrm{Coker}}}
\def\Coim{\mathop{\mathrm{Coim}}}

% Boxtimes
%
\DeclareMathSymbol{\boxtimes}{\mathbin}{AMSa}{"02}

%
% Macros for moduli stacks/spaces
%
\def\QCohstack{\mathcal{QC}\!\mathit{oh}}
\def\Cohstack{\mathcal{C}\!\mathit{oh}}
\def\Spacesstack{\mathcal{S}\!\mathit{paces}}
\def\Quotfunctor{\mathrm{Quot}}
\def\Hilbfunctor{\mathrm{Hilb}}
\def\Curvesstack{\mathcal{C}\!\mathit{urves}}
\def\Polarizedstack{\mathcal{P}\!\mathit{olarized}}
\def\Complexesstack{\mathcal{C}\!\mathit{omplexes}}
% \Pic is the operator that assigns to X its picard group, usage \Pic(X)
% \Picardstack_{X/B} denotes the Picard stack of X over B
% \Picardfunctor_{X/B} denotes the Picard functor of X over B
\def\Pic{\mathop{\mathrm{Pic}}\nolimits}
\def\Picardstack{\mathcal{P}\!\mathit{ic}}
\def\Picardfunctor{\mathrm{Pic}}
\def\Deformationcategory{\mathcal{D}\!\mathit{ef}}

% --- appended 2026-10-12 (Chief): the project own definitions of the reference/slogan/history
% display environments, copied mechanically from the shared Stacks preamble lines 8-12
% (\usepackage{verbatim} plus the three \newenvironment lines, which the chapter preamble keeps
% verbatim from the project preamble). Nothing added, removed or redefined: these environments
% carry no mathematical content and the project itself typesets them as comments. Without these
% three lines a counted statement whose own text uses them fails with Environment slogan undefined.
\usepackage{verbatim}
\newenvironment{reference}{\comment}{\endcomment}
\newenvironment{slogan}{\comment}{\endcomment}
\newenvironment{history}{\comment}{\endcomment}
\usepackage{cleveref}
\crefname{originallemma}{original lemma}{original lemmas}
\Crefname{originallemma}{Original lemma}{Original lemmas}
\makeatletter
\newlabel{algebra-lemma-finite-projective}{{00NX}{source}{}{sourcepointer-1}{}}
\newlabel{algebra-lemma-finite-projective@cref}{{[originallemma][0][]00NX}{[1][0][]source}{}{}{}}
\makeatother
\crefname{originaldefinition}{original definition}{original definitions}
\Crefname{originaldefinition}{Original definition}{Original definitions}
\makeatletter
\newlabel{definition-local-complete-intersection}{{8338}{source}{}{sourcepointer-2}{}}
\newlabel{definition-local-complete-intersection@cref}{{[originaldefinition][0][]8338}{[1][0][]source}{}{}{}}
\makeatother
\crefname{originallemma}{original lemma}{original lemmas}
\Crefname{originallemma}{Original lemma}{Original lemmas}
\makeatletter
\newlabel{lemma-base-change-NL}{{22771}{source}{}{sourcepointer-3}{}}
\newlabel{lemma-base-change-NL@cref}{{[originallemma][0][]22771}{[1][0][]source}{}{}{}}
\makeatother
\crefname{originallemma}{original lemma}{original lemmas}
\Crefname{originallemma}{Original lemma}{Original lemmas}
\makeatletter
\newlabel{lemma-cone-tor-amplitude}{{16972}{source}{}{sourcepointer-4}{}}
\newlabel{lemma-cone-tor-amplitude@cref}{{[originallemma][0][]16972}{[1][0][]source}{}{}{}}
\makeatother
\crefname{originallemma}{original lemma}{original lemmas}
\Crefname{originallemma}{Original lemma}{Original lemmas}
\makeatletter
\newlabel{lemma-lci-NL}{{22856}{source}{}{sourcepointer-5}{}}
\newlabel{lemma-lci-NL@cref}{{[originallemma][0][]22856}{[1][0][]source}{}{}{}}
\makeatother
\crefname{originallemma}{original lemma}{original lemmas}
\Crefname{originallemma}{Original lemma}{Original lemmas}
\makeatletter
\newlabel{lemma-n-pseudo-module}{{16003}{source}{}{sourcepointer-6}{}}
\newlabel{lemma-n-pseudo-module@cref}{{[originallemma][0][]16003}{[1][0][]source}{}{}{}}
\makeatother
\crefname{originallemma}{original lemma}{original lemmas}
\Crefname{originallemma}{Original lemma}{Original lemmas}
\makeatletter
\newlabel{lemma-pull-perfect}{{19017}{source}{}{sourcepointer-7}{}}
\newlabel{lemma-pull-perfect@cref}{{[originallemma][0][]19017}{[1][0][]source}{}{}{}}
\makeatother
\crefname{originallemma}{original lemma}{original lemmas}
\Crefname{originallemma}{Original lemma}{Original lemmas}
\makeatletter
\newlabel{lemma-pull-tor-amplitude}{{17148}{source}{}{sourcepointer-8}{}}
\newlabel{lemma-pull-tor-amplitude@cref}{{[originallemma][0][]17148}{[1][0][]source}{}{}{}}
\makeatother
\crefname{originallemma}{original lemma}{original lemmas}
\Crefname{originallemma}{Original lemma}{Original lemmas}
\makeatletter
\newlabel{lemma-tensor-NL}{{22754}{source}{}{sourcepointer-9}{}}
\newlabel{lemma-tensor-NL@cref}{{[originallemma][0][]22754}{[1][0][]source}{}{}{}}
\makeatother
\crefname{originallemma}{original lemma}{original lemmas}
\Crefname{originallemma}{Original lemma}{Original lemmas}
\makeatletter
\newlabel{lemma-tor-dimension}{{16997}{source}{}{sourcepointer-10}{}}
\newlabel{lemma-tor-dimension@cref}{{[originallemma][0][]16997}{[1][0][]source}{}{}{}}
\makeatother
\crefname{originallemma}{original lemma}{original lemmas}
\Crefname{originallemma}{Original lemma}{Original lemmas}
\makeatletter
\newlabel{lemma-two-out-of-three-perfect}{{18949}{source}{}{sourcepointer-11}{}}
\newlabel{lemma-two-out-of-three-perfect@cref}{{[originallemma][0][]18949}{[1][0][]source}{}{}{}}
\makeatother
\title{For a cocartesian square of local complete intersection maps the kernel of the cotangent comparison is finite projective}
\author{The Stacks Project Authors}
\date{stacks-project:more-algebra.tex:local-snapshot-20261012; GNU FDL 1.2 or later}
\begin{document}
\maketitle
\setcounter{subsection}{1}
\noindent\textit{Editorial note.} All delivered mathematical text is the author's own, copied from the pinned Stacks Project chapter file more-algebra.tex; no statement or proof step was rewritten, reordered or removed.\par\medskip
\noindent\textit{Editorial note.} The chapter file pinned-original-source.tex carries only \textbackslash{}input\{preamble\} on its first line, so no preamble text lives inside the chapter. The retained file sources/stacks-more-algebra/retained-preamble.tex is a mechanical, line-for-line copy of the project's single shared preamble as staged in the snapshot; nothing was added, removed or redefined and no package was added.\par\medskip
\noindent\textit{Editorial note.} The source URL is the source-level Stacks Project URL, as required when more than one card is produced from one source; the per-claim tag is carried by claim\_key.\par\medskip
\noindent\textit{Editorial note.} The counted statement span contains the author's cocartesian diagram, delivered verbatim. The eleven references inside the counted proof span are recorded as non-load-bearing navigation pointers; they mix cross-chapter tags with same-chapter label lines, and the applications that consume them stay inside the counted proof span.\par\medskip
\noindent\textit{Allowed standard prerequisites (editorial).}\begin{itemize}
\item The cotangent complex of a local complete intersection map being perfect of tor-amplitude in [-1, 0]. invoked at proof lines 22896-22897. Is the structural input identifying the two complexes as perfect with bounded amplitude.
\item Permanence of perfectness and of tor-amplitude under base change and the behaviour of these properties under distinguished triangles. invoked at proof lines 22898-22908. Is what makes the cone C perfect with the stated amplitude.
\item The base change computation for the cotangent complex, identifying the cohomology of the cone. invoked at proof lines 22908-22910. Locates the single nonzero cohomology module of C.
\end{itemize}\medskip
\ifdefined\hypertarget\hypertarget{sourcepointer-1}{}\fi
\noindent\textit{Source reference.} Original lemma 00NX: focused-seeds/AGP/sources/stacks-snapshot-20261011/tags.txt, source lines 875--875. Cross-chapter navigation pointer into the Stacks Project; the standard fact it supplies is applied inside the counted proof span and the argument resting on it stays there.\par\medskip
\ifdefined\hypertarget\hypertarget{sourcepointer-2}{}\fi
\noindent\textit{Source reference.} Original definition 8338: sources/stacks-more-algebra/pinned-original-source.tex, source lines 8338--8338. Same-chapter navigation pointer; the statement it names is standard for the section and the argument that consumes it stays inside the counted proof span.\par\medskip
\ifdefined\hypertarget\hypertarget{sourcepointer-3}{}\fi
\noindent\textit{Source reference.} Original lemma 22771: sources/stacks-more-algebra/pinned-original-source.tex, source lines 22771--22771. Same-chapter navigation pointer; the statement it names is standard for the section and the argument that consumes it stays inside the counted proof span.\par\medskip
\ifdefined\hypertarget\hypertarget{sourcepointer-4}{}\fi
\noindent\textit{Source reference.} Original lemma 16972: sources/stacks-more-algebra/pinned-original-source.tex, source lines 16972--16972. Same-chapter navigation pointer; the statement it names is standard for the section and the argument that consumes it stays inside the counted proof span.\par\medskip
\ifdefined\hypertarget\hypertarget{sourcepointer-5}{}\fi
\noindent\textit{Source reference.} Original lemma 22856: sources/stacks-more-algebra/pinned-original-source.tex, source lines 22856--22856. Same-chapter navigation pointer; the statement it names is standard for the section and the argument that consumes it stays inside the counted proof span.\par\medskip
\ifdefined\hypertarget\hypertarget{sourcepointer-6}{}\fi
\noindent\textit{Source reference.} Original lemma 16003: sources/stacks-more-algebra/pinned-original-source.tex, source lines 16003--16003. Same-chapter navigation pointer; the statement it names is standard for the section and the argument that consumes it stays inside the counted proof span.\par\medskip
\ifdefined\hypertarget\hypertarget{sourcepointer-7}{}\fi
\noindent\textit{Source reference.} Original lemma 19017: sources/stacks-more-algebra/pinned-original-source.tex, source lines 19017--19017. Same-chapter navigation pointer; the statement it names is standard for the section and the argument that consumes it stays inside the counted proof span.\par\medskip
\ifdefined\hypertarget\hypertarget{sourcepointer-8}{}\fi
\noindent\textit{Source reference.} Original lemma 17148: sources/stacks-more-algebra/pinned-original-source.tex, source lines 17148--17148. Same-chapter navigation pointer; the statement it names is standard for the section and the argument that consumes it stays inside the counted proof span.\par\medskip
\ifdefined\hypertarget\hypertarget{sourcepointer-9}{}\fi
\noindent\textit{Source reference.} Original lemma 22754: sources/stacks-more-algebra/pinned-original-source.tex, source lines 22754--22754. Same-chapter navigation pointer; the statement it names is standard for the section and the argument that consumes it stays inside the counted proof span.\par\medskip
\ifdefined\hypertarget\hypertarget{sourcepointer-10}{}\fi
\noindent\textit{Source reference.} Original lemma 16997: sources/stacks-more-algebra/pinned-original-source.tex, source lines 16997--16997. Same-chapter navigation pointer; the statement it names is standard for the section and the argument that consumes it stays inside the counted proof span.\par\medskip
\ifdefined\hypertarget\hypertarget{sourcepointer-11}{}\fi
\noindent\textit{Source reference.} Original lemma 18949: sources/stacks-more-algebra/pinned-original-source.tex, source lines 18949--18949. Same-chapter navigation pointer; the statement it names is standard for the section and the argument that consumes it stays inside the counted proof span.\par\medskip
\begin{lemma}
\label{lemma-base-change-lci-bis}
Consider a cocartesian diagram of rings
$$
\xymatrix{
B \ar[r] & B' \\
A \ar[r] \ar[u] & A' \ar[u]
}
$$
If $A \to B$ and $A' \to B'$ are local complete intersections as in
Definition \ref{definition-local-complete-intersection}, then
the kernel of $H^{-1}(\NL_{B/A} \otimes_B B') \to H^{-1}(\NL_{B'/A'})$
is a finite projective $B'$-module.
\end{lemma}

\begin{proof}
By Lemma \ref{lemma-lci-NL} the complexes $\NL_{B/A}$ and $\NL_{B'/A'}$
are perfect of tor-amplitude in $[-1, 0]$.
Combining Lemmas \ref{lemma-tensor-NL}, \ref{lemma-pull-perfect}, and
\ref{lemma-pull-tor-amplitude} we have
$\NL_{B/A} \otimes_B B' = \NL_{B/A} \otimes_B^\mathbf{L} B'$
and this complex is also perfect of tor-amplitude in $[-1, 0]$.
Choose a distinguished triangle
$$
C \to \NL_{B/A} \otimes_B B'  \to \NL_{B'/A'} \to C[1]
$$
in $D(B')$. By Lemmas \ref{lemma-two-out-of-three-perfect} and
\ref{lemma-cone-tor-amplitude} we conclude that $C$ is perfect
with tor-amplitude in $[-1, 1]$. By Lemma \ref{lemma-base-change-NL}
the complex $C$ has only one nonzero cohomology module, namely the module
of the lemma sitting in degree $-1$. This module is of finite presentation
(Lemma \ref{lemma-n-pseudo-module}) and flat
(Lemma \ref{lemma-tor-dimension}). Hence it is finite projective by
Algebra, Lemma \ref{algebra-lemma-finite-projective}.
\end{proof}

\end{document}
